\documentclass[11pt,a4paper]{article}

\usepackage[spanish,es-noshorthands]{babel}
\usepackage[T1]{fontenc}
\usepackage[utf8]{inputenc}
\usepackage{lmodern}
\usepackage[margin=2.4cm]{geometry}
\usepackage{microtype}
\usepackage{amsmath}
\usepackage{booktabs}
\usepackage{longtable}
\usepackage{array}
\usepackage{listings}
\usepackage{xcolor}
\usepackage{hyperref}
\usepackage{tikz}
\usetikzlibrary{arrows.meta,positioning,fit,backgrounds}
\usepackage{tcolorbox}
\usepackage{enumitem}
\usepackage{caption}
\usepackage{fancyhdr}
\setlength{\emergencystretch}{3em}

\hypersetup{
  colorlinks=true,
  linkcolor=black,
  urlcolor=blue!60!black,
  pdftitle={Manual de drs\_muones},
  pdfauthor={Laboratorio 5 - muones}
}

\pagestyle{fancy}
\fancyhf{}
\fancyhead[L]{\small drs\_muones: adquisición continua DRS4}
\fancyhead[R]{\small Laboratorio 5}
\fancyfoot[C]{\thepage}

\definecolor{codebg}{RGB}{248,248,248}
\definecolor{codekw}{RGB}{0,102,153}
\definecolor{codecm}{RGB}{90,120,90}

\lstdefinestyle{cpp}{
  language=C++,
  backgroundcolor=\color{codebg},
  basicstyle=\ttfamily\small,
  keywordstyle=\color{codekw}\bfseries,
  commentstyle=\color{codecm}\itshape,
  stringstyle=\color{red!60!black},
  numbers=none,
  frame=single,
  rulecolor=\color{gray!40},
  breaklines=true,
  columns=fullflexible,
  keepspaces=true,
  showstringspaces=false,
  tabsize=3
}

\lstdefinestyle{shell}{
  language=,
  backgroundcolor=\color{codebg},
  basicstyle=\ttfamily\small,
  frame=single,
  rulecolor=\color{gray!40},
  breaklines=true,
  columns=fullflexible,
  keepspaces=true,
  showstringspaces=false
}

\lstdefinestyle{py}{
  language=Python,
  backgroundcolor=\color{codebg},
  basicstyle=\ttfamily\small,
  keywordstyle=\color{codekw}\bfseries,
  commentstyle=\color{codecm}\itshape,
  stringstyle=\color{red!60!black},
  frame=single,
  rulecolor=\color{gray!40},
  breaklines=true,
  columns=fullflexible,
  keepspaces=true,
  showstringspaces=false
}

\newtcolorbox{aviso}[1][]{
  colback=orange!6,
  colframe=orange!60!black,
  boxrule=0.4pt,
  left=4pt,right=4pt,top=3pt,bottom=3pt,
  fonttitle=\bfseries,
  title=Advertencia,
  #1
}

\newtcolorbox{nota}[1][]{
  colback=blue!4,
  colframe=blue!40!black,
  boxrule=0.4pt,
  left=4pt,right=4pt,top=3pt,bottom=3pt,
  fonttitle=\bfseries,
  title=Nota,
  #1
}

\newcommand{\code}[1]{\texttt{#1}}

\title{\textbf{drs\_muones}\\[0.3em]
\large Adquisición continua con la DRS4 Evaluation Board\\
\large y almacenamiento binario}
\author{Laboratorio 5 -- experiencia de muones}
\date{Noviembre 2026}

\begin{document}

\maketitle

\begin{abstract}
\noindent
Este manual describe \code{drs\_muones}, un programa de adquisición continua para la
DRS4 Evaluation Board. El programa usa la librería del repositorio oficial
\code{drs4eb}: lee los cuatro canales de entrada, usa la entrada CH1 como trigger por
hardware, opcionalmente filtra los eventos (coincidencia entre canales y veto), los
encola en memoria RAM y los escribe en disco en un formato binario compacto y fácil
de leer desde Python. Está pensado para corridas largas (horas o días) con tasas
bajas, como la detección de muones cósmicos. Toda la configuración está en un solo
archivo, \code{src/config.h}.
\end{abstract}

\tableofcontents
\newpage

%=====================================================================
\section{Guía rápida}
%=====================================================================

\begin{enumerate}[itemsep=2pt]
  \item Abrir \code{nuevo\_codigo/src/config.h} con un editor de texto, cambiar los
        valores que hagan falta y guardar.
  \item En una terminal, parado en \code{nuevo\_codigo}, correr \code{./medir.sh}.
        El script compila (solo si cambió algo) y empieza a medir.
  \item Para terminar, apretar Ctrl+C \textbf{una sola vez} y esperar el resumen.
\end{enumerate}

Los datos quedan en la carpeta indicada en \code{CARPETA\_SALIDA}. Para probar sin
placa: \code{./medir.sh --selftest 100}.

\subsection{Cambios de noviembre de 2026}

\begin{itemize}[itemsep=2pt]
  \item \textbf{Posición del pulso.} Antes se configuraba \code{TRIGGER\_DELAY\_NS},
        que la placa cuenta hacia atrás desde el \emph{final} de la ventana: con
        600\,ns el pulso aparecía cerca de 1100\,ns. Ahora se configura
        \code{POSICION\_PULSO\_NS}, contado desde el \emph{comienzo} de la ventana
        (sección~\ref{sec:posicion}).
  \item \textbf{Carpeta de salida.} Si la carpeta se escribía sin barra final
        (\code{"ang\_0"}), los archivos quedaban afuera, con nombre
        \code{ang\_0run\_...bin}. Corregido: siempre quedan adentro de la carpeta.
  \item \textbf{Filtro de eventos} opcional: canales requeridos en coincidencia y
        canales de veto (sección~\ref{sec:filtro}).
  \item \textbf{Configuración} en castellano, con controles: un valor sin sentido da
        un mensaje claro al compilar. Ya no se pasan parámetros por CMake.
  \item \textbf{Archivo \code{\_info.txt}} por corrida, con la configuración y el
        resumen (triggers, guardados, descartados, duración).
  \item La línea de estado se imprime también mientras se espera un trigger.
  \item El formato binario no cambió de tamaño ni de orden: los scripts y los
        archivos anteriores siguen sirviendo (sección~\ref{sec:formato}).
\end{itemize}

%=====================================================================
\section{Hardware y conexionado}
%=====================================================================

La placa expone cuatro entradas SMA. Cada entrada está conectada a \emph{dos} canales
del chip DRS4, que se usan en pares para el modo de cascada:

\begin{center}
\begin{tabular}{ccc}
\toprule
Entrada SMA & Canal DRS usado & Columna en el archivo \\
\midrule
SMA1 & 0 & ch1 \\
SMA2 & 2 & ch2 \\
SMA3 & 4 & ch3 \\
SMA4 & 6 & ch4 \\
\bottomrule
\end{tabular}
\end{center}

El programa transfiere los canales 0 a 6 en una sola operación
(\code{TransferWaves(0, 6)}) y decodifica los cuatro canales pares. La columna de
tiempo \code{t} se toma de CH1.

\begin{nota}
El trigger es \textbf{por hardware en CH1} (entrada SMA1). Los otros canales no
disparan: se registran y, si el filtro está activado, se usan para decidir si el
evento se guarda.
\end{nota}

La entrada es de 1\,Vpp ($-0{,}5$\,V a $+0{,}5$\,V). Conviene atenuar la señal del PMT
o SiPM para no recortar; el programa marca los eventos saturados con un bit.

%=====================================================================
\section{Instalación y permisos}
%=====================================================================

\subsection{Dependencias}

En Debian/Raspberry Pi OS y en Ubuntu:

\begin{lstlisting}[style=shell]
sudo apt install build-essential cmake libusb-1.0-0-dev
\end{lstlisting}

Para los scripts de análisis: \code{sudo apt install python3-numpy python3-matplotlib}.

\subsection{Permisos de USB}

Sin una regla de udev, la placa solo es accesible como \code{root}. La regla está en
\code{codigo\_rp5/41-drs.rules}; el contenido es:

\begin{lstlisting}[style=shell]
ATTRS{idVendor}=="04b4", ATTRS{idProduct}=="1175", SUBSYSTEMS=="usb",
ACTION=="add", MODE="0666", GROUP="plugdev"
\end{lstlisting}

Para instalarla:

\begin{lstlisting}[style=shell]
sudo cp 41-drs.rules /etc/udev/rules.d/
sudo udevadm control --reload-rules
sudo udevadm trigger
\end{lstlisting}

y asegurarse de que el usuario pertenece al grupo \code{plugdev} (\code{groups}).
Tras reconectar la placa, queda accesible sin \code{sudo}.

%=====================================================================
\section{Compilación}
%=====================================================================

Lo normal es no compilar a mano: \code{./medir.sh} lo hace solo cada vez que cambia
\code{config.h}. Si falla, muestra la línea de \code{config.h} que tiene el problema
y deja el detalle en \code{build/compilacion.log}.

A mano, es:

\begin{lstlisting}[style=shell]
cd sobre_drs/nuevo_codigo
cmake -S . -B build
cmake --build build -j"$(nproc)"
\end{lstlisting}

El ejecutable queda en \code{build/drs\_muones}. La variable \code{DRS4EB\_DIR} indica
dónde está \code{software/} del repositorio oficial; por defecto apunta a
\code{../drs4eb/software}. Se compila optimizado (\code{Release}).

\begin{aviso}
Los parámetros son valores de compilación: cambiar \code{config.h} no tiene efecto
hasta volver a compilar. Con \code{./medir.sh} eso es automático; si se corre
\code{./build/drs\_muones} directamente, hay que compilar antes. Además, la carpeta
de salida es relativa al lugar desde donde se corre el programa: \code{./medir.sh}
siempre la crea dentro de \code{nuevo\_codigo}.
\end{aviso}

Las versiones anteriores permitían pisar parámetros con \code{cmake -DNOMBRE=valor}.
Eso se quitó para que haya un solo lugar donde mirar la configuración.

%=====================================================================
\section{Configuración: \code{config.h} en detalle}
%=====================================================================

Reglas para editar el archivo: cambiar solo el valor que está a la derecha del
nombre; decimales con punto; textos entre comillas dobles.

\begin{center}
\small
\begin{longtable}{p{4.6cm}p{2.0cm}p{8.0cm}}
\toprule
Parámetro & Defecto & Efecto \\
\midrule
\endhead
\code{CARPETA\_SALIDA} & \code{"medidas"} & Carpeta donde se guardan los archivos. Se crea si no existe. \\
\code{MINUTOS\_POR\_ARCHIVO} & 15 & Cada cuánto se cierra el archivo y se abre otro. \\
\code{FRECUENCIA\_GHZ} & 0.6 & Frecuencia de muestreo (0{,}1 a 6). Fija el paso temporal y la ventana. \\
\code{UMBRAL\_TRIGGER\_MV} & $-125$ & Nivel de trigger sobre CH1, en mV. \\
\code{PULSOS\_NEGATIVOS} & 1 & 1 = pulsos negativos, 0 = positivos. Vale para el trigger y el filtro. \\
\code{POSICION\_PULSO\_NS} & 300 & Dónde aparece el pulso de trigger, en ns desde el comienzo de la ventana. \\
\code{FILTRO\_ACTIVADO} & 0 & 0 = se guarda todo; 1 = se guarda solo lo que cumple el filtro. \\
\code{FILTRO\_CH1} \ldots\ \code{FILTRO\_CH4} & R, R, I, V & \code{REQUERIDO}, \code{IGNORAR} o \code{VETO}, por canal. \\
\code{UMBRAL\_CH1\_MV} \ldots\ \code{UMBRAL\_CH4\_MV} & $-125$ & Umbral de cada canal para decidir si tuvo un pulso. \\
\code{VENTANA\_COINCIDENCIA\_NS} & 50 & Separación máxima entre los pulsos de los canales \code{REQUERIDO}. \\
\code{SEGUNDOS\_ENTRE\_ESTADOS} & 5 & Periodicidad de la línea de estado. \\
\code{IMPRIMIR\_CADA\_EVENTO} & 0 & 1 = una línea por trigger, con la decisión del filtro. \\
\code{SATURACION\_MV} & 480 & Umbral para marcar un evento como saturado (no descarta). \\
\code{EVENTOS\_EN\_COLA} & 2000 & Tamaño de la cola en RAM ($\approx$ 20\,KB por evento). \\
\code{FORZAR\_ESCRITURA\_A\_DISCO} & 1 & \code{fsync} al cerrar cada archivo. \\
\code{ESPERA\_SONDEO\_US} & 500 & Espera entre consultas a la placa mientras se espera trigger. \\
\code{NS\_POR\_TICK} & 5.54 & Calibración de la posición del pulso (sección~\ref{sec:posicion}). \\
\bottomrule
\end{longtable}
\end{center}

Si un valor no tiene sentido (un modo de filtro mal escrito, un umbral con el signo
contrario al de los pulsos, una frecuencia fuera de rango), la compilación se
detiene con un mensaje en castellano que nombra el parámetro.

%---------------------------------------------------------------------
\subsection{Frecuencia de muestreo}
%---------------------------------------------------------------------

\code{FRECUENCIA\_GHZ} fija cuántas muestras por segundo toma el chip. El datasheet
garantiza de 0{,}7 a 5\,GSPS. Como la memoria analógica tiene 1024 celdas, la ventana
es
\[
  T_{\text{ventana}} = \frac{1024}{f_s}.
\]
A 0{,}6\,GSPS la ventana es de 1{,}70\,µs y el paso temporal 1{,}66\,ns; a 5\,GSPS la
ventana baja a 205\,ns y el paso a 0{,}2\,ns. La placa no da exactamente la
frecuencia pedida (pidiendo 0{,}6 da 0{,}6024); la real se imprime al arrancar y es
la que queda en el archivo.

\begin{aviso}
La calibración temporal de la placa vale para una sola frecuencia. Si no coincide con
la de medición, el programa avisa al arrancar y los tiempos guardados suponen
muestras de igual duración; el error llega a $\sim$10\,ns a 0{,}6\,GHz. Ver
\code{scripts/correccion\_temporal.html}.
\end{aviso}

%---------------------------------------------------------------------
\subsection{Umbral y polaridad del trigger}
%---------------------------------------------------------------------

\code{UMBRAL\_TRIGGER\_MV} es el nivel que debe cruzar CH1 para disparar.
\code{PULSOS\_NEGATIVOS} elige el flanco: 1 dispara cuando la señal baja del umbral,
0 cuando sube. El umbral es \textbf{absoluto}, no relativo a la línea de base. El
programa no compila si el signo del umbral no es coherente con la polaridad.

%---------------------------------------------------------------------
\subsection{Posición del pulso en la ventana}
\label{sec:posicion}
%---------------------------------------------------------------------

\code{POSICION\_PULSO\_NS} dice en qué tiempo de la ventana aparece el pulso que dio
el trigger, contado desde el comienzo (el borde izquierdo del gráfico). Con 300 a
0{,}6\,GHz, el pulso de CH1 aparece en $t \approx 300$\,ns: quedan 300\,ns de línea
de base antes y 1400\,ns después.

\subsubsection*{Por qué antes parecía estar al revés}

El chip muestrea sin parar, en círculo. Cuando llega el trigger, la placa espera un
tiempo y recién ahí frena; lo que se lee son los 1024 bins \emph{anteriores} al
frenado. Ese tiempo de espera es lo que la librería llama ``trigger delay''. Por eso
el retardo se cuenta hacia atrás desde el final de la ventana:
\[
  t_{\text{pulso}} = T_{\text{ventana}} - t_{\text{fijo}} - n_{\text{ticks}} \cdot \tau .
\]
$n_{\text{ticks}}$ (0 a 255) es lo que se programa en la placa, $\tau$ es lo que dura
un tick y $t_{\text{fijo}} = 23{,}5 + 28{,}2/f_s$\,ns es un retardo propio de la
placa (fórmula de la librería; 70\,ns a 0{,}6\,GHz). La versión anterior del programa
pasaba \code{TRIGGER\_DELAY\_NS} directo a la placa, y de ahí que con \code{delay=600}
el pulso apareciera en $\approx$1100\,ns y no en 600. El programa ahora despeja
$n_{\text{ticks}}$ de esa ecuación.

\subsubsection*{Calibración}

La librería supone $\tau = 6{,}2$\,ns. Medido con esta placa (serie 2566, firmware
21305) a 0{,}6024\,GHz, con las corridas de octubre y noviembre de 2026:

\begin{center}
\begin{tabular}{rrrr}
\toprule
\code{delay} pedido (ns) & Ticks & Pulso medido (ns) & Fórmula con $\tau=5{,}54$ (ns) \\
\midrule
100 & 16 & 1541 & 1541 \\
500 & 81 & 1184 & 1181 \\
600 & 97 & 1100 & 1092 \\
1000 & 161 & 734 & 738 \\
1740 & 255 (tope) & 220 & 217 \\
\bottomrule
\end{tabular}
\end{center}

El ajuste da $\tau = 5{,}54$\,ns por tick, con diferencias de hasta 8\,ns. Ese valor
está en \code{NS\_POR\_TICK}. Conviene contar con que el pulso quede a $\pm$10\,ns de
lo pedido. A otras frecuencias no hay mediciones: se usa la misma fórmula y hay que
verificar con datos dónde quedó el pulso.

\subsubsection*{Rango alcanzable}

Como hay solo 256 ticks, no se puede poner el pulso en cualquier lado:

\begin{center}
\begin{tabular}{rrr}
\toprule
$f_s$ (GHz) & Ventana (ns) & Posición alcanzable (ns) \\
\midrule
0.6 & 1700 & 220 a 1630 \\
1.0 & 1024 & 0 a 970 \\
2.0 & 512 & 0 a 475 \\
5.0 & 205 & 0 a 176 \\
\bottomrule
\end{tabular}
\end{center}

Si se pide un valor fuera del rango, el programa usa el más cercano y lo avisa. Si
se pide una posición mayor que la ventana, no arranca. Al arrancar siempre imprime
la posición que va a usar, y queda en el comentario del archivo (\code{pos=...ns}).

\begin{nota}
Para verificar o recalibrar: medir unos minutos, abrir un archivo con
\code{scripts/01\_ver\_un\_evento.py} y mirar en qué $t$ cruza CH1 el umbral. Si la
diferencia con lo pedido es sistemática, corregir \code{NS\_POR\_TICK}: un valor más
grande mueve el pulso hacia la derecha.
\end{nota}

%---------------------------------------------------------------------
\subsection{Filtro de eventos}
\label{sec:filtro}
%---------------------------------------------------------------------

Con \code{FILTRO\_ACTIVADO 0} se guarda todo evento que dispara CH1. Con
\code{FILTRO\_ACTIVADO 1}, cada evento se revisa antes de guardarlo. A cada canal se
le asigna uno de tres papeles:

\begin{center}
\begin{tabular}{lp{11cm}}
\toprule
Valor & Significado \\
\midrule
\code{REQUERIDO} & Tiene que haber un pulso en este canal. Si no, el evento se descarta. \\
\code{IGNORAR} & El canal no participa de la decisión (igual se guarda). \\
\code{VETO} & Si hay un pulso en este canal, el evento se descarta. \\
\bottomrule
\end{tabular}
\end{center}

``Hay un pulso'' significa que alguna muestra del canal pasó su umbral
(\code{UMBRAL\_CHn\_MV}), sin contar las 8 muestras de cada borde de la ventana,
donde el chip suele dejar picos espurios. El tiempo del pulso es el del primer cruce
del umbral, interpolado entre dos muestras.

El orden de la decisión es:

\begin{enumerate}[itemsep=2pt]
  \item Si algún canal \code{VETO} tiene pulso, en cualquier lugar de la ventana:
        descartado.
  \item Si a algún canal \code{REQUERIDO} le falta el pulso: descartado.
  \item Si entre el primer y el último pulso de los canales \code{REQUERIDO} hay más
        de \code{VENTANA\_COINCIDENCIA\_NS}: descartado.
  \item Si no, se guarda.
\end{enumerate}

La configuración que viene escrita (CH1 y CH2 requeridos, CH3 ignorado, CH4 veto,
ventana de 50\,ns) guarda los eventos con pulso casi simultáneo en CH1 y CH2 y sin
nada en CH4. Solo hace falta poner \code{FILTRO\_ACTIVADO} en 1.

\begin{aviso}
Lo que el filtro descarta \textbf{no se guarda} y no se puede recuperar. Antes de una
corrida larga conviene medir un rato sin filtro y probar el criterio sobre esos
datos con \code{scripts/02\_filtrar\_eventos.py}, que aplica la misma lógica. Como
referencia, en la corrida \code{eficiencia/132} (1446 eventos), pedir CH1 y CH2 con
$-125$\,mV deja 171 eventos, el 12\,\%: en la mayoría, el pulso de CH2 no llega a
$-125$\,mV.
\end{aviso}

Elección de los umbrales: en los datos de 2026 el ruido de CH1 a CH3 tiene picos de
entre $-40$ y $-65$\,mV sin que haya pulso, así que un umbral menos exigente que eso
da coincidencias falsas. Para un canal de veto sin nada conectado (CH4 quedó en
$-6 \pm 1$\,mV) cualquier umbral lejos de cero sirve.

La ventana de coincidencia: en los mismos datos, la diferencia $t_2 - t_1$ entre los
pulsos de CH2 y CH1 de un mismo muon está entre $-5$ y $+10$\,ns; 50\,ns deja margen
para el error de los tiempos sin calibrar ($\sim$10\,ns). El histograma que hace
\code{02\_filtrar\_eventos.py} sirve para elegirla.

Con \code{IMPRIMIR\_CADA\_EVENTO 1} se ve la decisión de cada trigger:

\begin{lstlisting}[style=shell]
trigger #2: guardado | CH1 en 303.0 ns CH2 en 304.6 ns CH3 en 306.5 ns CH4 sin pulso
trigger #3: descartado: falta el pulso en un canal REQUERIDO | CH1 en 301.9 ns ...
\end{lstlisting}

El campo \code{event\_index} de cada evento cuenta todos los triggers, también los
descartados, así que en el archivo los índices saltan. Los totales quedan en el
\code{\_info.txt}.

%---------------------------------------------------------------------
\subsection{Marca de saturación}
%---------------------------------------------------------------------

\code{SATURACION\_MV} no filtra: si alguna muestra de algún canal supera ese valor
en valor absoluto, el evento se guarda con el bit 0 de \code{flags} encendido.

%---------------------------------------------------------------------
\subsection{Cola en RAM y rotación de archivos}
%---------------------------------------------------------------------

\code{EVENTOS\_EN\_COLA} define cuántos eventos pueden esperar en memoria mientras el
hilo escritor los vuelca a disco. Cada evento ocupa 20\,496 bytes; con el valor por
defecto son unos 41\,MB. Si la cola se llena, la adquisición \textbf{espera} hasta
que haya lugar: no se pierde ningún evento ya leído, pero aumenta el tiempo muerto.

\code{MINUTOS\_POR\_ARCHIVO} controla cada cuánto se cierra el archivo y se abre uno
nuevo. A 0{,}25 triggers por segundo y 15 minutos, cada archivo pesa unos 4{,}6\,MB.
El cambio de archivo se hace al escribir un evento, así que un archivo puede durar
algo más si en ese momento no hay eventos.

%=====================================================================
\section{Cómo funciona por dentro}
\label{sec:paralelo}
%=====================================================================

El programa tiene dos hilos que trabajan en paralelo:

\begin{center}
\begin{tikzpicture}[node distance=7mm, every node/.style={font=\small},
  caja/.style={draw, rounded corners, inner sep=4pt, align=center}]
  \node[caja] (trig) {esperar\\trigger};
  \node[caja, right=of trig] (leer) {leer los\\4 canales};
  \node[caja, right=of leer] (filt) {filtro};
  \node[caja, right=of filt, fill=gray!15] (cola) {cola\\en RAM};
  \node[caja, right=of cola] (disco) {escribir\\al disco};
  \draw[-Stealth] (trig) -- (leer);
  \draw[-Stealth] (leer) -- (filt);
  \draw[-Stealth] (filt) -- (cola);
  \draw[-Stealth] (cola) -- (disco);
  \draw[-Stealth] (filt.south) -- ++(0,-5mm) -| (trig.south);
  \node[below=9mm of leer] {hilo principal};
  \node[below=9mm of disco] {hilo escritor};
\end{tikzpicture}
\end{center}

El hilo principal arma la placa, espera el trigger, trae las ondas por USB, aplica el
filtro y, si el evento pasa, lo copia a la cola. El hilo escritor saca eventos de la
cola y los escribe. Así una escritura lenta (una tarjeta SD, un disco ocupado) no
frena la medición mientras haya lugar en la cola.

Lo que \textbf{no} es paralelo es la placa: mientras se leen las ondas de un evento
(del orden de milisegundos por USB; no se midió en esta placa), no puede tomar otro. Ese tiempo muerto es propio del
hardware. A las tasas de esta experiencia (menos de un trigger por segundo) es
despreciable: menos del 1\,\% del tiempo.

El filtro corre en el hilo principal porque es mucho más rápido que la lectura USB
(recorre 4096 números) y así los eventos descartados ni siquiera ocupan lugar en la
cola.

%=====================================================================
\section{Formato binario}
\label{sec:formato}
%=====================================================================

Todos los valores son \textbf{little-endian}. El archivo empieza con un encabezado de
128 bytes y continúa con eventos de tamaño fijo. Hay una explicación para quien
nunca leyó un binario en \code{scripts/LEEME.md}.

\subsection{Encabezado de archivo (128 bytes)}

\begin{center}
\small
\begin{tabular}{rrl>{\raggedright\arraybackslash}p{7.4cm}}
\toprule
Offset & Tipo & Campo & Descripción \\
\midrule
0 & \code{char[8]} & \code{magic} & \code{"MUONBIN1"} \\
8 & \code{u16} & \code{version} & 2 (1 en los archivos anteriores al 14/11/2026) \\
10 & \code{u16} & \code{n\_channels} & 4 \\
12 & \code{u32} & \code{n\_samples} & 1024 \\
16 & \code{f64} & \code{sampling\_hz} & frecuencia real de muestreo \\
24 & \code{u16} & \code{board\_serial} & número de serie (0 en selftest) \\
26 & \code{u16} & \code{fw\_version} & firmware (0 en selftest) \\
28 & \code{u32} & \code{flags} & bit 0: \code{t} calibrado por la placa; bit 1: filtro activado \\
32 & \code{u64} & \code{start\_unix\_ns} & apertura del archivo, en ns desde 1970 \\
40 & \code{char[64]} & \code{comment} & configuración en texto \\
104 & \code{u8[24]} & \code{reserved} & ceros \\
\bottomrule
\end{tabular}
\end{center}

El comentario tiene la forma

\begin{lstlisting}[style=shell]
thr=-125mV pos=301ns freq=0.600GHz negEdge=1 filt=RRIV/50ns
\end{lstlisting}

donde \code{pos} es la posición del pulso que se usó y \code{filt} es \code{no} o
las letras de los cuatro canales (R, I, V) y la ventana. Los archivos de la versión
1 dicen \code{delay=} en lugar de \code{pos=}: ese número es el retardo que se le
pasaba a la placa, contado desde el final de la ventana.

\subsection{Evento}

\begin{center}
\small
\begin{tabular}{rrl>{\raggedright\arraybackslash}p{7.4cm}}
\toprule
Offset & Tipo & Campo & Descripción \\
\midrule
0 & \code{u64} & \code{unix\_ns} & instante en que se detectó el trigger \\
8 & \code{u32} & \code{event\_index} & número de trigger en la corrida (cuenta también los descartados) \\
12 & \code{u16} & \code{trigger\_cell} & celda de trigger \\
14 & \code{u16} & \code{flags} & bit 0: saturado; bits 1 a 4: pulso en CH1 a CH4 \\
16 & \code{f32[1024][5]} & datos & \{t[ns], ch1, ch2, ch3, ch4\} por muestra, en mV \\
\bottomrule
\end{tabular}
\end{center}

El tamaño de un evento es $16 + 1024 \times 5 \times 4 = 20\,496$ bytes y el del
archivo, para $N$ eventos, $128 + N \times 20\,496$.

\begin{nota}
La versión 2 tiene el mismo tamaño y orden que la 1. Lo único nuevo son bits que
antes valían cero (los bits 1 a 4 del evento y los \code{flags} del archivo). Los
bits de pulso se calculan con los umbrales \code{UMBRAL\_CHn\_MV} aunque el filtro
esté desactivado.
\end{nota}

\subsection{Archivo \code{\_info.txt}}

Junto a los \code{.bin} de cada corrida queda un archivo de texto con el mismo
prefijo. Se escribe al arrancar con la configuración completa y, al terminar bien,
se le agrega el resumen:

\begin{lstlisting}[style=shell]
Resumen de la corrida run_20261114_133529:
  duracion             : 3600.2 s (1.00 horas)
  triggers de CH1      : 912 (0.253 por seg)
  eventos guardados    : 118 (0.033 por seg)
  descartados          : 794
      falta un pulso REQUERIDO   : 781
      fuera de la ventana        : 9
      pulso en un canal VETO     : 4
  guardados y saturados: 0
  archivos .bin        : 5, en /home/.../nuevo_codigo/medidas
  escritura            : OK, los 118 eventos estan en disco
\end{lstlisting}

(Los números de este ejemplo son ilustrativos.) Para calcular tasas o eficiencias
hacen falta la duración y el total de triggers, que no se pueden sacar de los
\code{.bin} si hubo filtro.

%=====================================================================
\section{Uso}
%=====================================================================

\subsection{Modos}

\begin{lstlisting}[style=shell]
./medir.sh                  # compila si hace falta y mide hasta Ctrl+C
./medir.sh --selftest N     # N triggers inventados, sin placa
./build/drs_muones --help
\end{lstlisting}

Al arrancar se imprime la configuración efectiva:

\begin{lstlisting}[style=shell]
Placa DRS4 serie #2566, firmware 21305

Configuracion (se cambia en src/config.h):
  carpeta de salida : /home/.../nuevo_codigo/medidas
  archivo nuevo cada: 15 min
  frecuencia        : 0.6024 GHz (pedida 0.600) -> ventana de 1700 ns, una muestra cada 1.66 ns
  trigger           : CH1, umbral -125 mV, pulsos negativos
  posicion del pulso: 300 ns desde el comienzo de la ventana (pedida 300; alcanzable 217 a 1630; 240 ticks)
  filtro            : desactivado (se guarda todo lo que dispara CH1)
  cola en RAM       : 2000 eventos (41 MB)

Midiendo. Para terminar: Ctrl+C (una sola vez; espera a que guarde todo).

[12:31:05] triggers 342 (0.25 por seg) | guardados 342 | descartados 0 | cola 0/2000 | run_20261114_121500_idx0.bin
\end{lstlisting}

\subsection{Modo selftest}

\code{--selftest N} no usa la placa: genera $N$ eventos con pulsos inventados y los
pasa por el mismo filtro y la misma escritura. CH1 siempre tiene pulso, en la
posición configurada; CH2 lo tiene en coincidencia el 60\,\% de las veces y muy
tarde el 10\,\%; CH3 el 50\,\%; CH4 el 5\,\%. Sirve para probar un cambio de
configuración y para tener archivos con los que practicar los scripts. La secuencia
es siempre la misma.

\subsection{Cómo cortar una corrida}

\begin{center}
\begin{tabular}{lll}
\toprule
Acción & Señal & Comportamiento \\
\midrule
Ctrl+C & SIGINT & guarda lo que queda en la cola y cierra el archivo \\
\code{kill PID} & SIGTERM & ídem \\
corte de SSH & SIGHUP & ídem \\
segundo corte (después de $\geq$ 1\,s) & -- & salida inmediata, puede perder eventos \\
\bottomrule
\end{tabular}
\end{center}

El resumen final dice si todos los eventos guardados llegaron al disco.

\begin{aviso}
Para corridas largas por SSH, usar \code{tmux} para que la medición no dependa de la
conexión:
\begin{lstlisting}[style=shell]
tmux new -s muones
./medir.sh
# Ctrl+B, D para desacoplar; tmux attach -t muones para volver
\end{lstlisting}
\end{aviso}

\subsection{Qué mirar para saber que está midiendo bien}

\begin{itemize}[itemsep=2pt]
  \item \code{triggers} crece a un ritmo coherente con la tasa esperada.
  \item \code{cola} se mantiene cerca de cero: el disco sigue el ritmo.
  \item Con filtro, \code{guardados} crece; si queda en cero, el filtro descarta
        todo (probar con \code{IMPRIMIR\_CADA\_EVENTO 1}).
\end{itemize}

%=====================================================================
\section{Lectura de datos en Python}
%=====================================================================

Los ejemplos están en \code{scripts/}, y \code{scripts/LEEME.md} explica el formato
desde cero. Son scripts planos: arriba tienen las líneas a editar y se corren con
\code{python3 nombre.py} parado en \code{scripts/}.

\begin{center}
\begin{tabular}{lp{9.5cm}}
\toprule
Script & Qué hace \\
\midrule
\code{01\_ver\_un\_evento.py} & Lee el encabezado, lee los eventos y grafica uno. Comentado paso a paso. \\
\code{02\_filtrar\_eventos.py} & Aplica el criterio del filtro a una carpeta de archivos ya medidos, cuenta cuántos pasan y, si se pide, los guarda en archivos nuevos. \\
\code{03\_exportar\_a\_texto.py} & Pasa un evento a texto con columnas. \\
\code{04\_tasa\_por\_hora.py} & Tasa de eventos según la hora del día y a lo largo de la medición, a partir de la hora guardada en cada evento. \\
\code{read\_bin.py} & Corrección temporal por archivo (ver \code{correccion\_temporal.html}). \\
\bottomrule
\end{tabular}
\end{center}

Lo esencial de la lectura:

\begin{lstlisting}[style=py]
molde = np.dtype([('hora', '<u8'), ('numero', '<u4'), ('celda', '<u2'),
                  ('marcas', '<u2'), ('datos', '<f4', (1024, 5))])
eventos = np.fromfile(archivo, dtype=molde, offset=128)

t   = eventos[0]['datos'][:, 0]     # tiempos del primer evento, en ns
ch1 = eventos[0]['datos'][:, 1]     # CH1 del primer evento, en mV
\end{lstlisting}

Con todos los eventos a la vez:

\begin{lstlisting}[style=py]
horas     = eventos['hora']                      # ns desde 1970
minimos   = eventos['datos'][:, :, 1].min(axis=1)  # minimo de CH1 en cada evento
saturados = eventos['marcas'] & 1                # distinto de 0 si saturo
con_ch2   = eventos['marcas'] & 4                # distinto de 0 si CH2 tuvo pulso
\end{lstlisting}

El campo \code{hora} son los nanosegundos desde el 1/1/1970 (UTC), según el reloj de
la computadora en el momento en que detectó el trigger. Se pasa a fecha con
\code{datetime.datetime.fromtimestamp(hora / 1e9)}. Si esa computadora no tiene la
hora bien puesta (una Raspberry Pi sin internet la pierde al apagarse), las horas
guardadas están mal: conviene revisarla con \code{date} antes de medir.

Para archivos muy grandes, \code{np.memmap(archivo, dtype=molde, mode='r',
offset=128)} hace lo mismo sin cargar todo en memoria.

%=====================================================================
\section{Problemas frecuentes}
%=====================================================================

\begin{description}[style=nextline,itemsep=4pt]
  \item[\code{ERROR: no se pudo compilar}]
    Casi siempre es un valor mal escrito en \code{config.h}. El mensaje muestra la
    línea. Errores típicos: coma en lugar de punto decimal, una comilla que falta,
    \code{REQUERIDO} mal escrito.

  \item[\code{No se encontro ninguna placa DRS4.}]
    La placa no está conectada, no tiene alimentación USB suficiente
    ($\sim$500\,mA) o faltan permisos. Probar \code{lsusb} y revisar la regla de
    udev. Para probar sin hardware, \code{--selftest}.

  \item[\code{ERROR: POSICION\_PULSO\_NS ... no entra en la ventana}]
    Se pidió una posición más grande que la duración de la ventana a esa
    frecuencia. El mensaje indica el rango válido.

  \item[\code{triggers} no aumenta]
    No llegan pulsos que crucen el umbral. Revisar \code{UMBRAL\_TRIGGER\_MV},
    \code{PULSOS\_NEGATIVOS} y el cableado.

  \item[\code{triggers} aumenta pero \code{guardados} no]
    El filtro descarta todo. Con \code{IMPRIMIR\_CADA\_EVENTO 1} se ve el motivo
    de cada descarte.

  \item[Los archivos no aparecen donde se esperaba]
    La carpeta de salida es relativa al lugar desde donde se corre. Al arrancar se
    imprime la ruta completa.

  \item[La cola está siempre llena]
    El disco no sigue la tasa de triggers: activar el filtro, subir el umbral o
    usar un disco más rápido.

  \item[El último archivo quedó incompleto]
    El proceso murió sin poder cerrar (\code{kill -9}, corte de luz) o el disco se
    llenó. Los eventos completos que haya se pueden leer igual; la cantidad es la
    parte entera de (tamaño $-$ 128)/20\,496.

  \item[\code{Wrong firmware version}]
    Aviso de la librería con firmwares viejos. No impide medir.
\end{description}

%=====================================================================
\appendix
\section{Nombres anteriores de los parámetros}
%=====================================================================

\begin{center}
\small
\begin{tabular}{ll}
\toprule
Antes (octubre 2026) & Ahora \\
\midrule
\code{OUT\_DIR} & \code{CARPETA\_SALIDA} \\
\code{ROTATE\_MINUTES} & \code{MINUTOS\_POR\_ARCHIVO} \\
\code{SAMPLING\_GHZ} & \code{FRECUENCIA\_GHZ} \\
\code{THRESHOLD\_MV} & \code{UMBRAL\_TRIGGER\_MV} \\
\code{TRIGGER\_NEGATIVE\_EDGE} & \code{PULSOS\_NEGATIVOS} \\
\code{TRIGGER\_DELAY\_NS} & reemplazado por \code{POSICION\_PULSO\_NS} (otro significado) \\
\code{OVERLOAD\_MV} & \code{SATURACION\_MV} \\
\code{QUEUE\_EVENTS} & \code{EVENTOS\_EN\_COLA} \\
\code{FSYNC\_ON\_ROTATE} & \code{FORZAR\_ESCRITURA\_A\_DISCO} \\
\code{STATUS\_SECONDS} & \code{SEGUNDOS\_ENTRE\_ESTADOS} \\
\code{PRINT\_EVENTS} & \code{IMPRIMIR\_CADA\_EVENTO} \\
\code{POLL\_US} & \code{ESPERA\_SONDEO\_US} \\
\bottomrule
\end{tabular}
\end{center}

Para repetir una medición vieja hecha con \code{TRIGGER\_DELAY\_NS} $= D$ a
0{,}6\,GHz, la posición equivalente es aproximadamente $1630 - 0{,}89\,D$ (con el
tope de 220\,ns para $D \geq 1580$).

\section{Paso temporal y ventana según la frecuencia}

\begin{center}
\begin{tabular}{rrrr}
\toprule
$f_s$ (GSPS) & Paso (ns) & Ventana de 1024 bins (ns) & Bins en 100\,ns \\
\midrule
0.6 & 1.66 & 1700 & 60 \\
1.0 & 1.00 & 1024 & 100 \\
2.0 & 0.50 & 512 & 200 \\
5.0 & 0.20 & 205 & 500 \\
\bottomrule
\end{tabular}
\end{center}

\section{Referencias}

\begin{itemize}[itemsep=2pt]
  \item Repositorio oficial: \url{https://gitea.psi.ch/WaveDAQ/drs4eb} (librería
        \code{DRS.cpp}/\code{DRS.h}, ejemplo \code{drs\_exam.cpp}).
  \item \code{drs4eb/docs/manual/manual\_rev50.pdf}: manual de la placa V5.
  \item \code{DRS4\_rev09\_2.pdf}: datasheet del chip DRS4.
  \item \code{NOTAS\_DRS4\_muones.md}: notas del laboratorio.
  \item \code{scripts/correccion\_temporal.html}: corrección de tiempos por archivo.
\end{itemize}

\end{document}
